\cotefile{145}
\scanpage{20}
\mspage{8}

Il revient au $=$ de dire que $\sigma_0,\sigma_1,\sigma_\infty,\rho$ sont des générateurs, avec relations de définition

\[
(\overbrace{\rho^{-1}\sigma_\infty\sigma_0\deleted{\unsure{\rho^{-1}|\rho}}}^{l_\infty})
(\overbrace{\rho^{-1}\sigma_1\sigma_\infty\deleted{\unsure{\rho^{-1}|\rho}}}^{l_1})
(\overbrace{\rho^{-1}\sigma_0\sigma_1\deleted{\unsure{\rho^{-1}|\rho}}}^{l_0})=1 \tag*{(4)}
\]

\[
\sigma_0^2=\sigma_1^2=\sigma_\infty^2=1 \tag*{(5)}
\]

\[
\rho^3=1 \tag*{(6)}
\]

% slanted
\margin{\uline{NB} $\sigma_0(l_0)=l_{\deleted{\unsure{0}}\infty}^{-1}l_1^{-1}$\\
$\sigma_0(l_1)=l_\infty$\\
$\sigma_0(l_\infty)=l_1\deleted{\unsure{l_0^{-1}}}$\\
et itou pour $\sigma_1,\sigma_\infty$}

% slanted
\margin{Il faut ajouter la relation\unsure{,} $\rho\sigma_0\rho^{-1}=\sigma_1$ \deleted{etc} On a deux gén. $\sigma_0,\rho$ en fait avec relations $\sigma_0^2=\rho^3=1$, la \deleted{$[\unsure{\underbrace{\rho^{-1}\sigma_0\rho}_{\sigma_\infty}\sigma_0\rho^{-1}\underbrace{\rho\sigma_0\rho^{-1}}_{\sigma_1}\underbrace{\rho^{-1}\sigma_0\rho}_{\sigma_\infty}\rho^{-1}}]$} relation (4) étant conséquence des \unsure{précédentes|précédtes} (quand on y fait $\sigma_1=\rho\sigma_0\rho^{-1}$, $\sigma_\infty=\rho^{-1}\sigma_0\rho$). Noter que\\
$l_0=\deleted{\unsure{(\sigma_0\rho)^2}}\rho^{-1}(\sigma_0\rho)^2\rho=\boxed{\unsure{\rho^{-1}\sigma_0\rho\sigma_0\rho^{-1}}}$\\
$l_1=(\deleted{\unsure{\rho}}\sigma_0\rho)^2$\\
$l_\infty=(\rho^{-1}\sigma_0\deleted{\unsure{\rho}})^2$\\
$\mathrm{int}(\sigma_0)\,l_1=l_\infty$\\
$\mathrm{int}(\sigma_0)\,l_\infty=l_1$\\
$\mathrm{int}(\sigma_0)(l_0)=l_\infty^{-1}l_1^{-1}$}

\uline{NB} Si on y fait \inserted{de plus} $\sigma_0,\sigma_1,\sigma_\infty=1$, on trouve comme relation de définition $\rho^3=1$, \deleted{\unsure{donc}} sans plus, ce qui prouve que $\rho$ n'est pas un mot p.r. aux $\sigma_i$ \ldots

Tout ceci est compatible avec les notations du début, \deleted{car dans ces dernières on avait tout pris au \unsure{premiers|premier}, de \unsure{signes|gén.} \unsure{intervertis|interposés} $l'_i=l_i^{-1}$, $\sigma'_i=\sigma_{\dot\rho(i)}$, \unsure{donc} \uline{\unsure{remplacer}} $\sigma'$} \deleted{\unsure{puis} $\sigma_i$} si on y remplace les $l_i$ par les $l_i^{-1}$. Ici on a, comme $\sigma_i=\sigma'_{\unsure{\dot\rho^{-1}(i)|\rho^{-1}(i)|\dot\rho(i)}}$

\[
\left\{\begin{aligned}
\mathrm{int}(\sigma_i)&:\ \text{échange } l_j \text{ et }l_{\deleted{\unsure{j}}k}
 \ \inserted{(i,j,k)\ \text{distincts}}\quad\deleted{\text{et}}\\[-2pt]
&\deleted{\text{(et sa valeur sur }l_{\unsure{\dot\rho^{-1}(i)|\rho^{-1}(i)|\dot\rho(i)}}\text{ s'en déduit par (1))}}\\
\mathrm{int}(\rho)\,l_i&=l_{\dot\rho(i)}
\end{aligned}\right. \tag*{(7)}
\]

On voudrait maintenant expliciter les $\varepsilon_i$ avec $\varepsilon_i^2=l_i$ $(i=0,1,\infty)$, $\varepsilon_i$ étant donc un générateur \unsure{du|des} centralisateur (ou normalisateur, c'est itou) des $\unsure{L_i|l_i|b_i}$ dans ${\unsure{E|\Sigma|\mathscr E}}^+$. (ou \unsure{le centralisateur|les centralisateurs} dans $E$). Les calculs faits plus haut, avec les chgts de notation indiqués, donnent

\[
\left\{\begin{aligned}
\varepsilon_0&=l_0^{\unsure{1|-1}}\sigma_1l_{\unsure{1}}^{\unsure{1|-1}}\\
\varepsilon_1&=l_\infty^{-1}\sigma_\infty l_1^{-1}\\
\varepsilon_\infty&=l_0^{-1}\sigma_0 l_\infty^{-1}
\end{aligned}\right. \tag*{\deleted{\unsure{(8)}}}
\]
% archive sheet number: 19
